\documentclass[a4paper,10pt]{article}

% ==============================================================================
% ATS-OPTIMIZED CV TEMPLATE
%
% Reference: VitaeContext - cv-ats module
% This LaTeX template is designed to pass Applicant Tracking Systems (ATS)
% by ensuring a strict top-to-bottom, left-to-right plain text extraction,
% while maintaining a clean, professional human-readable PDF layout.
%
% Formatting Rules referenced:
% - hub/cv-ats/formatting-rules.md (Single-column layout, standard fonts)
% - hub/cv-ats/agent-workflow.md (Macro-driven content injection)
% - hub/cv-ats/common-pitfalls.md (No hidden tables, use \hfill for alignment)
% ==============================================================================

% --- PACKAGES ---
\usepackage[utf8]{inputenc}
\usepackage[T1]{fontenc}
\usepackage[scaled]{helvet} % Standard system font (Helvetica/Arial clone)
\renewcommand{\familydefault}{\sfdefault}
\usepackage[left=1.4cm,right=1.4cm,top=1.2cm,bottom=1.2cm]{geometry}
\usepackage{enumitem} % For strict bullet point control without ATS breakage
\setlist[itemize]{label=\textbullet} % Force standard solid circle bullet
\setlist{leftmargin=*,noitemsep,topsep=0.05em,partopsep=0em,parsep=0em}
\usepackage[none]{hyphenat} % Disable auto-hyphenation to prevent keyword splitting
\usepackage{parskip}
\setlength{\parskip}{0.1em}
\setlength{\parindent}{0pt}
\setlength{\emergencystretch}{3em}
\hfuzz=2pt
\usepackage{xcolor}
\definecolor{myblue}{HTML}{205792} % Professional dark blue
\usepackage[colorlinks=true,urlcolor=myblue,linkcolor=myblue,citecolor=myblue]{hyperref}
\usepackage{microtype}

% --- MACROS ---
% The \cventry macro uses \hfill to align dates to the right.
% ATS parsers will read "Role Date" sequentially from left to right,
% avoiding the "invisible table hack" failure (see hub/cv-ats/common-pitfalls.md).
\newcommand{\cventry}[2]{\noindent\begin{minipage}[t]{\linewidth}\textbf{#1}\hfill\textbf{#2}\end{minipage}}

\hypersetup{pdftitle={FirstName LastName CV},pdfauthor={FirstName LastName}}
\pagestyle{empty}
\setcounter{secnumdepth}{0}

% --- SECTION TITLE STYLE ---
\newcommand{\mysection}[1]{%
  \vspace{0.3ex}
  {\bfseries\scshape\color{myblue} #1}\enspace%
  \leaders\hrule height 0.4pt\hfill\kern0pt
  \vspace{-0.2em}
}

\begin{document}

% ==============================================================================
% HEADER (CONTACT INFORMATION)
% Reference: hub/cv-ats/core-sections.md
% Rule: Place contact info in the document body, not in page headers/footers.
% Rule: Write out URLs explicitly (e.g., linkedin.com/in/username).
% ==============================================================================
\begin{center}
    {\LARGE\scshape\color{myblue} [FIRST NAME] [LAST NAME]}\\[0.15em]
    {\normalsize [Target Job Title] | [Key Specialization]}\\[0.25em]
    \href{tel:+1234567890}{+1 234 567 890} \quad
    \href{mailto:email@example.com}{email@example.com} \quad
    \href{https://linkedin.com/in/username}{linkedin.com/in/username} \quad
    \href{https://github.com/username}{github.com/username}\\
    {[City, State/Country]} \enspace|\enspace [Work Authorization] \enspace|\enspace \href{https://portfolio.com}{portfolio.com}
\end{center}

% ==============================================================================
% PROFESSIONAL SUMMARY (OPTIONAL BUT RECOMMENDED)
% Reference: hub/cv-ats/core-sections.md
% Rule: Use standard title "Professional Summary". Not "About Me".
% Reference: hub/cv-ats/keyword-strategy.md
% Rule: Integrate target job description keywords naturally.
% ==============================================================================
\mysection{PROFESSIONAL SUMMARY}

[A 3-4 sentence elevator pitch. Focus on your years of experience, core technical competencies, and your most significant quantifiable achievement. Mirror the exact phrasing of the job description for hard skills.]

% ==============================================================================
% SKILLS
% Reference: hub/cv-ats/keyword-strategy.md
% Rule: Differentiate between hard skills and soft skills. Group them logically.
% ==============================================================================
\mysection{SKILLS}
\begin{itemize}
    \item \textbf{Category 1 (e.g., Security Engineering):} [Skill 1], [Skill 2], [Skill 3]
    \item \textbf{Category 2 (e.g., Programming):} [Skill 1], [Skill 2], [Skill 3]
    \item \textbf{Category 3 (e.g., Systems \& Platforms):} [Skill 1], [Skill 2], [Skill 3]
\end{itemize}

% ==============================================================================
% WORK EXPERIENCE
% Reference: hub/cv-ats/core-sections.md
% Rule: Reverse-chronological order. Use one explicit date format throughout.
% Reference: hub/cv-ats/achievement-metrics.md
% Rule: Use STAR or XYZ formula. Action Verb + Task + Result (Metric).
% ==============================================================================
\mysection{WORK EXPERIENCE}

\cventry{[Job Title], [Company Name] - [Location]}{[MM/YYYY] - Present}
\begin{itemize}
    \item .[Action Verb] [Specific Task/Project] [Result/Impact with quantifiable metric \%, \$, \#]. (XYZ Method: Accomplished X as measured by Y, by doing Z).
    \item Developed [System/Feature] utilizing [Hard Skill/Tool], resulting in [Quantifiable Business Outcome].
    \item Led [Project/Initiative] by [Action Taken], achieving [Measurable Improvement].
\end{itemize}

\vspace{0.2em}
\cventry{[Previous Job Title], [Company Name] - [Location]}{[MM/YYYY] - [MM/YYYY]}
\begin{itemize}
    \item .[Action Verb] [Task] to achieve [Result]. Include exact keywords from the JD.
    \item Managed [Process/Team], improving [Metric] by [Number]\%.
\end{itemize}

% ==============================================================================
% KEY PROJECTS (OPTIONAL)
% Use this section if you lack professional experience or want to highlight
% open-source contributions.
% ==============================================================================
\mysection{KEY PROJECTS}
\vspace{0.2em}

\cventry{\textbf{[Project Name]} - \href{https://github.com/username/project}{github.com/username/project}}{[MM/YYYY]}
\begin{itemize}
    \item Built a [App Type] using [Tech Stack] that solves [Specific Problem].
    \item Achieved [Metric/Result] by implementing [Feature/Optimization].
    \item \textbf{Stack:} [Technology 1], [Technology 2], [Technology 3]
\end{itemize}

% ==============================================================================
% EDUCATION
% Reference: hub/cv-ats/core-sections.md
% Rule: Standard naming ("Education").
% ==============================================================================
\mysection{EDUCATION}
\cventry{[Degree Name (e.g., MSc in Cybersecurity)]\enspace|\enspace [University Name], [Country]}{[MM/YYYY] - [MM/YYYY]}\\
{[Optional: GPA, Honors, Relevant Coursework, or Thesis Title]}

% ==============================================================================
% CERTIFICATIONS & LANGUAGES
% ==============================================================================
\mysection{CERTIFICATIONS \& LANGUAGES}
\begin{itemize}
    \item \textbf{Certifications:} [Certification Name] ([Issuing Organization], [Year])
    \item \textbf{Languages:} [Language 1] ([Proficiency]), [Language 2] ([Proficiency])
\end{itemize}

\end{document}
